%% Section 4.3: delay and the approximant logical relation. First draft.
%% Budget: 2.75 pages (notes/paper-outline.md). The definitions follow the
%% Agda: agda/Mu/Compiler/Delay.agda (the pass),
%% agda/Mu/Semantics/Graph/Clos3.agda and Clos4.agda (the semantics),
%% agda/Mu/proof/Correctness/Graph/DelayRel.agda (the relation),
%% DelayReflect.agda (the fundamental lemma, collect),
%% DelayPreserve.agda (the forward relation) and DelayCorrect.agda (the
%% theorem). The history of the relation is told as a tour of relation
%% shapes, without names or dates (notes/history-outline.md).

\newcommand{\Val}{\mathsf{Value}}
\newcommand{\den}[1]{[\![#1]\!]}
\newcommand{\app}{\mathbin{\bullet}}
\newcommand{\callop}{\mathbin{\odot}}
\newcommand{\pairv}[2]{\langle\!|#1,#2|\!\rangle}
\newcommand{\car}{\mathsf{car}}
\newcommand{\cdr}{\mathsf{cdr}}
\newcommand{\ALR}{\mathsf{ALR}}
\newcommand{\Calls}{\mathsf{Calls}}
\newcommand{\depth}{\mathsf{depth}}
\newcommand{\kw}[1]{\mathtt{#1}}

\subsection{Delay and the Approximant Logical Relation}
\label{sec:approx-lr}

This section is about the one pass of our compiler whose proof is not
an equation. Delay is the pass of closure conversion that changes how a
closure is represented. In its source language, Clos3, a closure
$\kw{clos}\;N\;\vec{M}$ pairs the code $N$, which is closed and takes
its free variables as a tuple, with the expressions $\vec{M}$ that
compute those free variables; the closure is still a function, and an
application $\kw{app}\;L\;M$ applies it directly. In the target
language, Clos4, a closure is an ordinary pair, and a call takes the
pair apart:
\[
\begin{array}{rcl}
  \mathit{delay}(\kw{clos}\;N\;\vec{M}) & = &
    \kw{pack}\;(\kw{code}\;\mathit{delay}(N))\;(\kw{tuple}\;\mathit{delay}(\vec{M}))\\
  \mathit{delay}(\kw{app}\;L\;M) & = &
    \kw{callc}\;\mathit{delay}(L)\;\mathit{delay}(M)
\end{array}
\]
Every other construct is translated homomorphically.

\paragraph{The two semantics}
Recall from Section~\ref{sec:graph-models} that a denotation is a set
of values, that a function's values are entries $V \mapsto w$ with a
finite, nonempty list $V$ of inputs, and that application collects the
outputs of the entries whose inputs the argument covers:
\[
  D \app E \;=\; \{\, w \mid V \mapsto w \in D,\; V \neq [\,],\; V \subseteq E \,\}.
\]
In Clos3, a closure means its code applied to the tuple of its free
variables, so it denotes a set of function entries, and an application
is $\den{\kw{app}\;L\;M}\rho = \den{L}\rho \app \den{M}\rho$. In Clos4,
code is a curried function of the environment and then the argument,
$\kw{pack}$ builds the set of pairs $\pairv{u}{v}$ with $u$ from the
code and $v$ from the environment, and a call applies the first
components to the second components and then to the argument:
\[
  \den{\kw{callc}\;L\;M}\rho \;=\; (\car(\den{L}\rho) \app \cdr(\den{L}\rho)) \app \den{M}\rho,
\]
where $\car(D) = \{u \mid \pairv{u}{v} \in D\}$ and
$\cdr(D) = \{v \mid \pairv{u}{v} \in D\}$. We write $D' \callop E'$ for
the right-hand side with $\den{L}\rho = D'$ and $\den{M}\rho = E'$.

\paragraph{Why preservation is easy and reflection is not}
Delay changes the representation of values, so its correctness is not
an equation. The theorem we prove is that a closed program and its
translation produce the same constants, and that one produces something
exactly when the other does. The forward half, preservation, needs a
relation between a \emph{single} source value and a \emph{set} of
target values, defined by recursion on the source value. Its interesting
case says what it means for a target set to implement a function entry:
\[
  (U \mapsto w) \sqsubseteq D' \iff
  D' \text{ contains a pair, and }
  U \sqsubseteq^* E' \text{ implies } w \sqsubseteq D' \callop E'
  \text{ for every consistent } E'.
\]
The source value is finite, so the recursion is structural, and the
proof that $v \in \den{M}\rho$ implies
$v \sqsubseteq \den{\mathit{delay}(M)}\rho'$ is a routine induction on
$M$.

The backward half, reflection, must show that every value the target
produces is accounted for by the source, and here the representation of
closures gets in the way. To see how, consider a function that adds
nine, $\lambda y.\, x + y$ in an environment where $x$ is $9$. In
Clos3 its meaning is a set of entries such as $[1] \mapsto 10$ and
$[6] \mapsto 15$. After delay, its code takes the environment as an
extra argument, and the meaning of the code is a table with entries for
\emph{every} environment, including $[x{=}0] \mapsto ([6] \mapsto 6)$,
even though the program only ever pairs this code with the environment
$x{=}9$. So the target's meanings contain code--environment
combinations that correspond to no source function. Worse, a call takes
the code from $\car$ and the environment from $\cdr$
\emph{separately}, so a result may combine the code entry of one pair
with the environment of another. No single target value corresponds to
a single source value, and the relation that worked for preservation
has no counterpart for reflection.

\paragraph{Relations that do not work}
We, and others we consulted, tried the natural alternatives, and each
failure points to a piece of the solution.
\begin{itemize}
\item \emph{Relating single values}, a target value to a source value,
  works for every construct except the call, because of the mixing of
  $\car$ and $\cdr$ described above.
\item \emph{Relating whole sets on both sides} handles the mixing, but
  the function case then mentions the relation in both positive and
  negative positions, so the definition is not well founded. A
  step-indexed relation would guard the recursive occurrence with a
  ``later'' modality and discharge it at a step of evaluation; a
  denotational semantics has no steps.
\item \emph{Indexing by the depth of the related values}, in the style
  of the projections of the inverse-limit construction, fails because a
  shallow output of a function can depend on an arbitrarily deep input.
\item \emph{Indexing by one target value within its whole target set}
  comes closest, but the function case must relate the entire argument
  set, which is infinite and cannot be shrunk to match the finite index.
\end{itemize}

\paragraph{The key insight}
A simulation needs only one side to shrink. Reflection asks whether
what the target produces is produced by the source, so the target side
can be a \emph{finite} list of observations, while the source side
remains an arbitrary, possibly infinite, set. And because a call mixes
the $\car$ and $\cdr$ of several pairs, the matching of codes to
environments must happen \emph{inside} the finite list, so the index is
a list of values rather than a single value.

\paragraph{The relation}
For a finite list $V$ of target values, a \emph{batch of calls}
$\Calls(V, U, W)$ records that each result $w$ in the finite list $W$
comes from an entry
$\pairv{\mathit{FV} \mapsto (U_0 \mapsto w)}{t}$ of $V$ whose
environment inputs $\mathit{FV}$ are nonempty and are among the
environments $\cdr(V)$ that $V$ itself provides, and whose argument
inputs $U_0$ are nonempty and contained in $U$; and that each $u$ in
$U$ is an argument input of such an entry. The approximant logical
relation $\ALR(D, V)$ relates a source set $D$ to a finite list $V$ of
target values when:
\begin{enumerate}
\item every constant in $V$ is in $D$;
\item if $V$ is nonempty, so is $D$;
\item for every batch of calls $\Calls(V, U, W)$ and every consistent
  source set $E$ with $U$ consistent and $\ALR(E, U)$, we have
  $\ALR(D \app E, W)$;
\item for the projection $\pi$ of each tuple component or sum
  injection, and every finite $W \subseteq \pi(V)$, we have
  $\ALR(\pi(D), W)$.
\end{enumerate}
Clause~3 is the only place where $\car$ and $\cdr$ meet, and they meet
inside the finite list $V$. The definition is by recursion on the
depth of $V$: the arguments $U$ and the results $W$ of a batch of calls
are strictly shallower than $V$, and so are the projections in
clause~4. In Agda the recursion is expressed with a counter,
$\ALR_0 = \top$ and
$\ALR_{k+1}(D, V) = \mathcal{L}(\ALR_k)(D, V)$, where $\mathcal{L}(Q)$
is clauses 1--4 with the recursive occurrences replaced by $Q$, and
$\ALR(D, V) = \ALR_{1 + \depth(V)}(D, V)$. The counter is only a
termination device: any counter that bounds the depth of $V$ gives the
same relation (Lemma~\ref{lem:depth-irrelevant}), which lets us fold
and unfold $\ALR$ as if it were defined directly by clauses 1--4.

\begin{lemma}[Depth irrelevance]\label{lem:depth-irrelevant}
If every value in $V$ has depth less than $k$ and less than $j$, then
$\ALR_k(D, V)$ if and only if $\ALR_j(D, V)$.
\end{lemma}

The relation is also monotone in $D$ and antitone in $V$: a larger
source set and a sublist of the target list preserve it.

\paragraph{Why not recursion on types}
The intermediate languages are typed, so one might hope to define the
relation by recursion on types, as \citet{Chlipala:2007fk} and
\citet{Nielsen:2000aa} do. Recursive types rule this out. The meaning of
$\kw{fold}$ and $\kw{unfold}$ is the identity, so the relation at
$\mu X.\,A$ would have to be the relation at $A[\mu X.\,A/X]$, which is
not smaller; at $\mu X.\,X \Rightarrow \mathrm{Int}$, the application
clause asks for the relation at the argument type, which is
$\mu X.\,X \Rightarrow \mathrm{Int}$ itself. The depth of the finite
approximation, by contrast, decreases at every call, and the values of
the graph model are untyped, so the depth needs no help from the types.

\paragraph{The fundamental lemma}
Environments $\rho$ and $\rho'$ are related when, for every variable
$x$ and every finite $V \subseteq \rho'(x)$, we have
$\ALR(\rho(x), V)$, together with the consistency and nonemptiness
that the graph model requires of environments.

\begin{lemma}[Fundamental lemma]\label{lem:delay-fundamental}
If $\rho$ and $\rho'$ are related, then for every term $M$ of Clos3 and
every finite $V \subseteq \den{\mathit{delay}(M)}\rho'$, we have
$\ALR(\den{M}\rho, V)$.
\end{lemma}

The proof is by induction on $M$. Two cases carry the weight. For a
call, a finite $V \subseteq D' \callop E'$ is produced by finitely many
pairs and arguments, by continuity of the graph model: there are finite
$W \subseteq D'$ and $U \subseteq E'$ with $\Calls(W, U, V)$. The
induction hypotheses give $\ALR(\den{L}\rho, W)$ and
$\ALR(\den{M}\rho, U)$, and clause~3 of the first, applied to this batch
of calls, gives $\ALR(\den{L}\rho \app \den{M}\rho, V)$. For a closure,
we must establish clause~3 for the set of pairs
$\den{\kw{pack}\;(\kw{code}\;N')\;(\kw{tuple}\;\vec{M}')}\rho'$: an
entry of a batch of calls has environment inputs $\mathit{FV}$ drawn
from the environments that the list itself provides, which are
approximations of the tuple of free variables, so the induction
hypothesis for $\vec{M}$ relates them to the source environment, and
the induction hypothesis for the body $N$ relates the results. This is
where the condition $\mathit{FV} \subseteq \cdr(V)$ in the definition
of $\Calls$ earns its keep: it is exactly what excludes the code
entries for environments that the program never supplies.

\begin{theorem}[Correctness of delay]\label{thm:delay}
For a closed term $M$ of Clos3 and a constant $c$,
$c \in \den{M}$ if and only if $c \in \den{\mathit{delay}(M)}$, and
$\den{M}$ is nonempty if and only if $\den{\mathit{delay}(M)}$ is.
\end{theorem}

Reflection follows from Lemma~\ref{lem:delay-fundamental} with the
empty environments: a constant $c$ of the target is the one-element
list $[c]$, and clauses 1 and 2 give $c \in \den{M}$ and the
nonemptiness of $\den{M}$. Preservation follows from the single-value
relation $\sqsubseteq$.

\paragraph{The cost}
Delay is the most expensive pass of the denotational proof: 524 lines
of Agda, against 35 to 139 for the same pass in the three operational
styles (Section~\ref{sec:evaluation}). Most of the difference has three
sources. First, the denotational proof needs two relations, one for
each direction, while a simulation proof gets the backward direction
for free from forward simulation and the determinism of the target.
Second, reflection needs the finite lists, the batches of calls, and
the lemma that extracts them from a target call. Third, $\beta$-rules in
the graph model hold only under side conditions of nonemptiness and
continuity, which a machine gets for free by taking a step. We return
to these numbers in Section~\ref{sec:evaluation}.
